% GENERATED FILE. Edit canonical lessons, then run npm run generate:books.
% canonical-source-hash: fnv1a64:018df617635819a4
% canonical-lessons: MR-C463-garva, MR-C463-ahankar, MR-C463-svabhiman, MR-C463-saramne, MR-C463-osalne

\chapter{Pride and Shame}
\label{ch:mr-pride-and-shame}
\hlchaptermodality{463}

\begin{chapteropening}
\textbf{By the end of this chapter:} \emph{I can talk about pride, ego and feeling ashamed.}

Complete the last lesson of chapter 463: I can talk about pride, ego and feeling ashamed.
\end{chapteropening}

\section[\texorpdfstring{\mr{गर्व} (garva)}{गर्व (garva)}]{\mr{गर्व} (garva) — pride; conceit}

\label{lesson:MR-C463-garva}



\begin{quote}
Before the new one: say the Marathi for temptation, then the Marathi for greed.
\end{quote}

\subsection*{You'll want to know: \mr{गर्व}}
\textbf{\mr{गर्व}} — \emph{garva} — \textquotedblleft{}pride; conceit\textquotedblright{}.

\begin{quote}
\mr{बक्षीस} \mr{मिळाल्याचा} \mr{त्याला} \mr{गर्व} \mr{झाला}.
\end{quote}

\begin{quote}
\emph{bakṣīs miḷālyācā tyālā garva jhālā.}
\end{quote}

\begin{quote}
Winning the prize made him conceited.
\end{quote}

Masculine.

Often a fault: \mr{गर्व} \mr{होणे}, to grow conceited; \mr{गर्विष्ठ}, arrogant.

\begin{cousinweb}
\textbf{\mr{गर्व}} (\emph{garva}) is a tatsama, taken from Sanskrit \emph{garva}, \textquotedblleft{}pride, arrogance\textquotedblright{} (CDIAL 4064).
\end{cousinweb}

\begin{grammarlens}[title={What you've built}]
One more word for describing what you lived and felt.
\end{grammarlens}

\subsection*{Your turn}
\begin{itemize}
\raggedright
  \item \emph{Say it:} \emph{garva}
  \item \emph{Say it:} \emph{garva}, once more
  \item \emph{Say it:} \emph{bakṣīs miḷālyācā tyālā garva jhālā.} — Winning the prize made him conceited.
  \item \emph{Recall:} say the Marathi for temptation, then the Marathi for greed, then say \emph{garva} again
\end{itemize}

\begin{tcolorbox}[breakable,colback=teal!4,colframe=teal!35!black,title={Before you move on}]
What does \mr{गर्व} mean? (\textquotedblleft{}Pride; conceit\textquotedblright{}.) Say it once more.
\end{tcolorbox}
\section[\texorpdfstring{\mr{अहंकार} (ahaṅkār)}{अहंकार (ahaṅkār)}]{\mr{अहंकार} (ahaṅkār) — ego; arrogance}

\label{lesson:MR-C463-ahankar}



\begin{quote}
Before the new one: say the Marathi for greed, then the Marathi for pride.
\end{quote}

\subsection*{You'll want to know: \mr{अहंकार}}
\textbf{\mr{अहंकार}} — \emph{ahaṅkār} — \textquotedblleft{}ego; arrogance\textquotedblright{}.

\begin{quote}
\mr{त्याच्या} \mr{अहंकारामुळे} \mr{मित्र} \mr{दूर} \mr{गेले}.
\end{quote}

\begin{quote}
\emph{tyācyā ahaṅkārāmuḷe mitra dūr gele.}
\end{quote}

\begin{quote}
His friends drifted away because of his ego.
\end{quote}

Masculine.

\mr{अहं}, I, + \mr{कार}, making: the sense of I.

\begin{cousinweb}
\textbf{\mr{अहंकार}} (\emph{ahaṅkār}) is a tatsama, taken from Sanskrit \emph{ahaṃkāra}, \textquotedblleft{}conception of one's individuality\textquotedblright{} (Monier-Williams), built on \emph{aham}, \textquotedblleft{}I\textquotedblright{}.
\end{cousinweb}

\begin{grammarlens}[title={What you've built}]
One more word for describing what you lived and felt.
\end{grammarlens}

\subsection*{Your turn}
\begin{itemize}
\raggedright
  \item \emph{Say it:} \emph{ahaṅkār}
  \item \emph{Say it:} \emph{ahaṅkār}, once more
  \item \emph{Say it:} \emph{tyācyā ahaṅkārāmuḷe mitra dūr gele.} — His friends drifted away because of his ego.
  \item \emph{Recall:} say the Marathi for greed, then the Marathi for pride, then say \emph{ahaṅkār} again
\end{itemize}

\begin{tcolorbox}[breakable,colback=teal!4,colframe=teal!35!black,title={Before you move on}]
What does \mr{अहंकार} mean? (\textquotedblleft{}Ego; arrogance\textquotedblright{}.) Say it once more.
\end{tcolorbox}
\section[\texorpdfstring{\mr{स्वाभिमान} (svābhimān)}{स्वाभिमान (svābhimān)}]{\mr{स्वाभिमान} (svābhimān) — self-respect}

\label{lesson:MR-C463-svabhiman}



\begin{quote}
Before the new one: say the Marathi for pride, then the Marathi for ego.
\end{quote}

\subsection*{You'll want to know: \mr{स्वाभिमान}}
\textbf{\mr{स्वाभिमान}} — \emph{svābhimān} — \textquotedblleft{}self-respect\textquotedblright{}.

\begin{quote}
\mr{गरीब} \mr{असला} \mr{तरी} \mr{त्याने} \mr{स्वाभिमान} \mr{सोडला} \mr{नाही}.
\end{quote}

\begin{quote}
\emph{garīb aslā tarī tyāne svābhimān soḍlā nāhī.}
\end{quote}

\begin{quote}
Though he was poor, he never gave up his self-respect.
\end{quote}

Masculine.

\mr{स्व}, own, + \mr{अभिमान}, pride; \mr{स्वाभिमानी}, self-respecting.

\begin{cousinweb}
\textbf{\mr{स्वाभिमान}} (\emph{svābhimān}) is a learned compound built from two Sanskrit words, \emph{sva}, \textquotedblleft{}own\textquotedblright{}, and \emph{abhimāna}, \textquotedblleft{}pride\textquotedblright{} (Monier-Williams).
\end{cousinweb}

\begin{grammarlens}[title={What you've built}]
One more word for describing what you lived and felt.
\end{grammarlens}

\subsection*{Your turn}
\begin{itemize}
\raggedright
  \item \emph{Say it:} \emph{svābhimān}
  \item \emph{Say it:} \emph{svābhimān}, once more
  \item \emph{Say it:} \emph{garīb aslā tarī tyāne svābhimān soḍlā nāhī.} — Though he was poor, he never gave up his self-respect.
  \item \emph{Recall:} say the Marathi for pride, then the Marathi for ego, then say \emph{svābhimān} again
\end{itemize}

\begin{tcolorbox}[breakable,colback=teal!4,colframe=teal!35!black,title={Before you move on}]
What does \mr{स्वाभिमान} mean? (\textquotedblleft{}Self-respect\textquotedblright{}.) Say it once more.
\end{tcolorbox}
\section[\texorpdfstring{\mr{शरमणे} (śaramṇe)}{शरमणे (śaramṇe)}]{\mr{शरमणे} (śaramṇe) — to feel ashamed}

\label{lesson:MR-C463-saramne}



\begin{quote}
Before the new one: say the Marathi for ego, then the Marathi for self-respect.
\end{quote}

\subsection*{You'll want to know: \mr{शरमणे}}
\textbf{\mr{शरमणे}} — \emph{śaramṇe} — \textquotedblleft{}to feel ashamed\textquotedblright{}.

\begin{quote}
\mr{आपली} \mr{चूक} \mr{लक्षात} \mr{आल्यावर} \mr{तो} \mr{शरमला}.
\end{quote}

\begin{quote}
\emph{āplī cūk lakṣāt ālyāvar to śaramlā.}
\end{quote}

\begin{quote}
He felt ashamed when he realised his mistake.
\end{quote}

Past (he …): \mr{तो} \mr{शरमला}, \emph{to śaramlā}.

From \mr{शरम}, shame, a Persian word beside the Marathi \mr{लाज}.

\begin{cousinweb}
\textbf{\mr{शरमणे}} (\emph{śaramṇe}) is a Marathi verb built on \mr{शरम}, \textquotedblleft{}shame\textquotedblright{}, a loan from Persian \emph{sharm}.
\end{cousinweb}

\begin{grammarlens}[title={What you've built}]
One more word for describing what you lived and felt.
\end{grammarlens}

\subsection*{Your turn}
\begin{itemize}
\raggedright
  \item \emph{Say it:} \emph{śaramṇe}
  \item \emph{Say it:} \emph{śaramṇe}, once more
  \item \emph{Say it:} \emph{āplī cūk lakṣāt ālyāvar to śaramlā.} — He felt ashamed when he realised his mistake.
  \item \emph{Recall:} say the Marathi for ego, then the Marathi for self-respect, then say \emph{śaramṇe} again
\end{itemize}

\begin{tcolorbox}[breakable,colback=teal!4,colframe=teal!35!black,title={Before you move on}]
What does \mr{शरमणे} mean? (\textquotedblleft{}To feel ashamed\textquotedblright{}.) Say it once more.
\end{tcolorbox}
\section[\texorpdfstring{\mr{ओशाळणे} (ośāḷṇe)}{ओशाळणे (ośāḷṇe)}]{\mr{ओशाळणे} (ośāḷṇe) — to feel sheepish}

\label{lesson:MR-C463-osalne}



\begin{quote}
Before the new one: say the Marathi for self-respect, then the Marathi for to feel ashamed.
\end{quote}

\subsection*{You'll want to know: \mr{ओशाळणे}}
\textbf{\mr{ओशाळणे}} — \emph{ośāḷṇe} — \textquotedblleft{}to feel sheepish\textquotedblright{}.

\begin{quote}
\mr{खोटं} \mr{पकडलं} \mr{गेलं} \mr{तेव्हा} \mr{तो} \mr{ओशाळला}.
\end{quote}

\begin{quote}
\emph{khoṭaṁ pakaḍlaṁ gelaṁ tevhā to ośāḷlā.}
\end{quote}

\begin{quote}
He felt sheepish when his lie was caught out.
\end{quote}

Past (he …): \mr{तो} \mr{ओशाळला}, \emph{to ośāḷlā}.

Embarrassment after being caught out: \mr{ओशाळलेला} \mr{चेहरा}.

\begin{cousinweb}
\textbf{\mr{ओशाळणे}} (\emph{ośāḷṇe}) is a Marathi verb whose origin is not settled.
\end{cousinweb}

\begin{grammarlens}[title={What you've built}]
One more word for describing what you lived and felt.
\end{grammarlens}

\subsection*{Your turn}
\begin{itemize}
\raggedright
  \item \emph{Say it:} \emph{ośāḷṇe}
  \item \emph{Say it:} \emph{ośāḷṇe}, once more
  \item \emph{Say it:} \emph{khoṭaṁ pakaḍlaṁ gelaṁ tevhā to ośāḷlā.} — He felt sheepish when his lie was caught out.
  \item \emph{Recall:} say the Marathi for self-respect, then the Marathi for to feel ashamed, then say \emph{ośāḷṇe} again
\end{itemize}

\begin{tcolorbox}[breakable,colback=teal!4,colframe=teal!35!black,title={Before you move on}]
What does \mr{ओशाळणे} mean? (\textquotedblleft{}To feel sheepish\textquotedblright{}.) Say it once more.
\end{tcolorbox}
